\naslov{Ansambli napovednih modelov}

V ansamblu $\hat{M}$ imamo več modelov $\hat{m}_i$.
Ideja je, da vrednosti ciljnih spremenljivk kombiniramo iz napovedi posameznih
modelov v ansamblu.
Običajno to naredimo s povprečjem, za razvrščanje pa vzamemo gostiščnico.
Za razvrščanje lahko pridemo do pričakovane napake:
če vsakemu od $T$ modelov pripišemo naključno spremenljivko $Z_i = \mathbbm{1}(y
\ne \hat{y}_i)$, in predpostavimo, da so te spremenljivke neodvisne z $E(Z_i) =
\varepsilon_i$, potem je
\begin{align*}
  \operatorname{Err}(\hat{M})
  &= P\!\left( \sum_i Z_i > \frac{T}{2} \right)
  = P\!\left( \sum_i (Z_i - E(Z_i)) + T \varepsilon > \frac{T}{2} \right) \\
  &= P\!\left( \inv{T} \sum_i (Z_i - E(Z_i)) > \pol - \varepsilon \right)
  \le \exp\!\left( -2 T (\pol - \varepsilon)^2 \right)
\end{align*}
po Hoeffdingovi neenakosti.
Napaka je torej reda $e^{-T}$, če je le $\varepsilon < \pol$.

\vprasanje{Oceni napako klasifikacijskega ansambla.}

Obravnavajmo homogen ansambel regresijskih modelov.
Razlike potem pridejo iz različnih podatkovnih množic $S$, ki jih obravnavajmo
kot slučajno spremenljivko.
\[
  E_S(\hat{M}(\mb{x}_0))
  = E_S(\inv{T} \sum_i \hat{m}_i(\mb{x}_0))
  = \inv{T} \sum_i E_S(\hat{m}_i(\mb{x}_0))
  = E_S(\hat{m}(x_0))
\]
Torej je pristranskost $E_S(\hat{M}(\mb{x}_0)) - m(\mb{x}_0)$ enaka kot
pristranskost osnovnega modela $\hat{m}$.

\vprasanje{Pokaži, da ansambel ne spremeni pristranskosti.}

Izračunamo lahko, da za korelacijski faktor
\[
  \rho_S(\mb{x}_0)
  = \frac{E_S((\hat{m}_i(\mb{x}_0) - E(\hat{m}_i(\mb{x}_0))) (\hat{m}_j(\mb{x}_0) -
	E(\hat{m}_j(\mb{x}_0))))}{\sqrt{ \var_S(\hat{m}_i(\mb{x}_0))
	  \var_S(\hat{m}_j(\mb{x}_0)) }}
\]
velja
\[
  \var_S(\hat{M}(\mb{x}_0))
  = \left( \rho + \frac{1-\rho}{T} \right) \var_S(\hat{m}(\mb{x}_0)).
\]
Levemu členu produkta pravimo \pojem{faktor zmanjševanja variance}.

\vprasanje{Kaj je faktor zmanjševanja variance? Izpelji predpis.}

Za pridobivanje različnih modelov v homogenem ansamblu imamo več možnosti.
Ena od njih je vrečenje, kjer spreminjamo učne podatke za vsak model z
vzorčenjem skupne učne množice s ponavljanjem.
Napako ansambla lahko potem ocenimo kot povprečje napak na učni množici, kjer
napovedi za posamičen primer iz učne množice generiramo le z modeli, ki tega
testa niso imeli med učnimi podatki.
Temu pravimo \pojem{ocena out-of-bag}.

\vprasanje{Kaj je ocena out-of-bag?}

Drug način za pridobivanje različnih modelov so naključni podprostori, kjer vsak
model naučimo na $q < p$ naključno izbranih spremenljivkah z vsemi podatki.
Tu ni vzorčenja, tako da ne moramo uporabljati ocen out-of-bag.

% LocalWords:  gostiščnico Hoeffdingovi vrečenje generiramo out-of-bag
